\section{The \bench{} benchmark}
\label{sec:benchmark}

\bench{} is a deterministic, resettable sandbox of six simulated services, twelve parameterized task
templates, a fault injector that operates at the service boundary, and a grader that reads only the
ground-truth effect ledger. Table~\ref{tab:services} summarizes the services.

\begin{table}[t]
\centering
\small
\caption{Services and tool contracts. ``Key'' marks writes that accept an idempotency key in the native
contract; lags are documented in the tool descriptions.}
\label{tab:services}
\resizebox{\linewidth}{!}{%
\begin{tabular}{lll}
\toprule
Service & Writes (idempotency) & Read-back path (consistency) \\
\midrule
social & \texttt{publish} (key on mastodon only) & \texttt{list\_posts}: weibo lags 180\,s; none for x \\
billing & \texttt{create\_charge} (key), \texttt{refund} (naturally idempotent) & \texttt{list\_charges} (strong) \\
tickets & \texttt{create}, \texttt{add\_comment} (non-idempotent), \texttt{update\_status} (conditional) & \texttt{list\_recent} (strong), \texttt{search} (lags 120\,s) \\
mail & \texttt{send} (non-idempotent, irreversible) & \texttt{search\_sent} (lags 120\,s) \\
data & \texttt{insert}, \texttt{insert\_many} (non-atomic), \texttt{upsert} (idempotent) & \texttt{query} (strong) \\
deploy & \texttt{trigger} (non-idempotent) & \texttt{list\_runs}, \texttt{get\_run} (strong) \\
\bottomrule
\end{tabular}}
\end{table}

\paragraph{Services and contracts.} Semantics follow well-known production conventions. Billing keys follow
Stripe: a repeated key with the same parameters replays the original charge and a repeated key with different
parameters is rejected \citep{stripe_idempotent}. Social posts, tickets, emails, database rows and deployment
runs are created anew by every successful call unless a key is honored. Several read paths are eventually
consistent with a documented lag (a search index, a mail Sent folder, one social platform's listing), and
one platform has no listing endpoint, so its writes cannot be verified. Search endpoints behave like search
engines: a record matches when every query term occurs in its fields. Three general tools complete the
interface: \texttt{wait} advances a virtual clock; \texttt{escalate\_to\_human} reaches a simulated on-call
operator who inspects the ground truth of every write that returned an error and answers after 15 simulated
minutes (or, in one condition, never answers); and \texttt{finish} records a structured final report with a
status and a list of operations whose outcome the agent considers uncertain. Every tool additionally
carries a machine-readable contract in the style of MCP tool annotations \citep{mcp2025spec}: read-only,
idempotent and destructive hints, the fields that identify a write's intent, a read-back function, the
documented visibility lag and a compensating action. The agent never sees the contract object; only the
harness-level recovery policies in \S\ref{sec:setup} may read it.

\paragraph{Tasks.} Twelve templates describe realistic operational workflows: announcing a release on two
platforms, opening a ticket with the post links and emailing a summary; billing two or three customers for an
invoice and notifying finance; opening an incident and paging on-call; recording a batch of database
migrations; rolling out a service to staging and then production; refunding a duplicate charge while
keeping the original; enabling feature flags and announcing them; notifying three customers individually;
cross-posting an event; upgrading a subscription; resolving an incident ticket; and a long-horizon hotfix
rollout with about ten writes across five services. Each template is instantiated with seeded parameters
(products, versions, customers, amounts) and states its targets precisely (for example, \emph{exactly one}
post containing the release version on each named platform). Each instruction ends with one sentence asking for
every action to happen exactly once; E6 removes it to measure how much this cue drives caution. Each template also
declares its \emph{focal
writes}---the calls a fault may attach to---annotated with idempotency class (non-idempotent, key optional,
naturally idempotent, idempotent), verification class (strong, eventual, none) and reversibility. Across
templates there are 34--35 focal writes per instance.

\paragraph{Fault injection.} A fault is attached to the $n$-th call that matches a focal write's tool and
argument filter, so the same world and the same fault apply to every model, harness and condition. Table
\ref{tab:faults} lists the twelve modes. The observable response of each ambiguous mode is byte-identical
across its hidden outcome states, and the virtual clock advances identically, so a pair such as
\texttt{timeout\_pre}/\texttt{timeout\_post} differs only in whether the service committed. A request that
the service would reject anyway receives its real error and leaves the trigger armed for the next matching
call; this keeps the pairs symmetric even when an agent's first attempt is malformed. Late commits land
90 simulated seconds after the request was sent (\texttt{timeout\_late}) or after a delay drawn log-uniformly
between 40\,s and 2\,h (\texttt{timeout\_late\_tail}, median about 9 minutes), whatever the agent does in the
meantime; the delay is drawn deterministically per world, so every model and policy faces the same one. At the end
of an episode, requests still in flight complete, as they would in reality.

\begin{table}[t]
\centering
\small
\caption{Fault modes. Rows in the same group are observation-equivalent at the time of the response.}
\label{tab:faults}
\resizebox{\linewidth}{!}{%
\begin{tabular}{llll}
\toprule
Mode & Agent observes & Hidden outcome & Exactly-once recovery \\
\midrule
\texttt{timeout\_pre} & timeout & not executed & verify, then re-issue \\
\texttt{timeout\_post} & timeout & executed, response lost & verify, then skip \\
\texttt{timeout\_late} & timeout & executed 90\,s later (still in flight) & same key, or wait and escalate \\
\texttt{timeout\_late\_tail} & timeout & executed 40\,s--2\,h later (heavy tail) & same key \\
\midrule
\texttt{http500\_pre} & HTTP 500 & not executed & verify, then re-issue \\
\texttt{http500\_post} & HTTP 500 & executed & verify, then skip \\
\midrule
\texttt{partial\_timeout} & timeout on a batch & first half executed & verify, re-issue only missing rows \\
\texttt{duplicate\_delivery} & success & executed twice & send a key proactively \\
\midrule
\texttt{http503\_transient} & 503, retry later & not executed & retry \\
\texttt{rate\_limit} & 429 with retry-after & not executed & wait, then retry \\
\texttt{outage} & persistent 503 & not executed & stop and report \\
\texttt{schema\_drift} & 400, field renamed & not executed & repair arguments \\
\texttt{none} & success & executed & -- \\
\bottomrule
\end{tabular}}
\end{table}

\paragraph{Grading.} The grader reads the ledger of committed effects and the final world state; no model is
used as a judge. It reports task success (TS), exactly-once success (EOS), the number of duplicate
executions (including compensated ones), residual duplicates, collateral damage to records the task did not
target (for example refunding the original charge), unrequested writes, and \emph{overclaim}: finishing
with status \texttt{completed} although the task failed or a duplicate remains. From the agent's own tool
calls---not the harness's---a deterministic coder classifies the recovery behaviour that follows the focal
fault as blind retry, retry with the same key, retry with a new key, verify then retry, verify then skip,
escalate, stop without checking, or move on.

\paragraph{Validity checks.} Thirty-one unit tests with scripted agents check the fault semantics exactly:
observation equivalence of each pair, symmetric handling of invalid requests, stale reads before and fresh
reads after the documented lag, partial batches, key replay, conflict and batch resumption, late commits
defeating verification but not keys, the heavy-tailed delay distribution, the waiting thresholds of the
wait-then-verify policies, redelivery, escalation reporting ground truth, and the grader's detection of
collateral damage and overclaiming. Scripted reference policies behave as the propositions predict: a
blind-retry agent duplicates only in committed worlds, a verify-first agent is exactly-once except under late
commits, and a same-key agent is exactly-once under the transient faults and resumable-key contract tested here.
In the empirical data, agents cannot distinguish the
observation-equivalent worlds (\S\ref{sec:results}), confirming that no information about the hidden outcome
leaks through the response.

\paragraph{Harness-agnostic deployment.} The sandbox runs as a local HTTP service owned by the experiment
runner. A stdio MCP server forwards only \texttt{tools/list} and \texttt{tools/call} to it, so any
MCP-capable agent can be evaluated unchanged; the fault plan and the ground truth never leave the runner.
Each harness episode runs in an empty working directory with a throwaway harness home, so memory and session
state cannot carry over between episodes.
